\PassOptionsToPackage{table}{xcolor}
\documentclass[10pt,aspectratio=169]{beamer}
\newcommand{\LectureNumber}{05}
\input{course-theme.tex}
\title{Types and arithmetic expressions}
\subtitle{CSC103 Programming Fundamentals / Lecture 05}
\author{Dr. Muhammad Farhan}
\institute{Associate Professor of Computer Science\\Department of Computer Science\\COMSATS University Islamabad\\Sahiwal Campus}
\date{Fall 2026}
\begin{document}
\begin{frame}\titlepage\end{frame}

\begin{frame}[fragile,t]{The problem for this lecture}
\begin{columns}[T,onlytextwidth]\begin{column}{0.60\textwidth}
Convert 3672 seconds to hours, minutes and seconds using / and \%.
\end{column}\begin{column}{0.35\textwidth}\small
Find complete hours with seconds / 3600.
\end{column}\end{columns}
\note{Introduce the target problem before explaining the tools. Revisit it in the guided solution later.\par Syllabus reading: Liang Ch 2. CLO-1.}
\end{frame}

\begin{frame}[fragile,t]{Learning objectives}
\begin{columns}[T,onlytextwidth]\begin{column}{0.60\textwidth}
\begin{itemize}
\item Choose suitable primitive types
\item Evaluate precedence and integer division
\item Use casts and conversions deliberately
\end{itemize}
\end{column}\begin{column}{0.35\textwidth}\small
CLO-1

Primitive data types\par Division and remainder\par Precedence and grouping\par Conversions and numeric limits
\end{column}\end{columns}
\note{Explain how each objective supports the opening problem. The final questions check these objectives.\par Syllabus reading: Liang Ch 2. CLO-1.}
\end{frame}

\begin{frame}[fragile,t]{Primitive data types}
\begin{itemize}
\item Integer types: byte, short, int and long.
\item Floating-point types: float and double. Other primitives: char and boolean.
\item A String is a reference type.
\end{itemize}
\vspace{1mm}\begin{block}{Example}
\begin{lstlisting}
int count = 20;
long population = 8000000000L;
double average = 72.5;
boolean passed = true;
\end{lstlisting}
\end{block}
\note{byte/short/int/long use 8/16/32/64 bits. float/double use 32/64 bits. char is 16 bits. Do not claim a fixed Java language storage size for boolean.\par Syllabus reading: Liang Ch 2. CLO-1.}
\end{frame}

\begin{frame}[fragile,t]{Division and remainder}
\begin{itemize}
\item Integer division truncates toward zero.
\item The \% operator returns a remainder. Its sign follows the dividend.
\item At least one floating-point operand gives floating-point division.
\end{itemize}
\vspace{1mm}\begin{block}{Example}
\begin{lstlisting}
7 / 2       // 3
7 / 2.0     // 3.5
7 % 2       // 1
-7 / 2      // -3
\end{lstlisting}
\end{block}
\note{A later assignment to double does not undo integer division that already occurred. Division by zero with integral operands throws ArithmeticException.\par Syllabus reading: Liang Ch 2. CLO-1.}
\end{frame}

\begin{frame}[fragile,t]{Precedence and grouping}
\begin{itemize}
\item Multiplication, division and remainder bind before addition and subtraction.
\item Operators at these arithmetic levels associate left to right.
\item Parentheses make the intended calculation explicit.
\end{itemize}
\vspace{1mm}\begin{block}{Example}
\begin{lstlisting}
2 + 3 * 4       // 14
(2 + 3) * 4     // 20
20 / 5 * 2      // 8
\end{lstlisting}
\end{block}
\note{Walk through intermediate values. Avoid mixing side effects with arithmetic while students learn precedence.\par Syllabus reading: Liang Ch 2. CLO-1.}
\end{frame}

\begin{frame}[fragile,t]{Conversions and numeric limits}
\begin{itemize}
\item Widening conversions often happen automatically.
\item A narrowing cast can discard a fraction or lose range.
\item Floating-point values approximate many decimal fractions.
\end{itemize}
\vspace{1mm}\begin{block}{Example}
\begin{lstlisting}
double x = 7;
int y = (int) 7.9;  // 7
double mean = (double) 7 / 2;
\end{lstlisting}
\end{block}
\note{Some widening conversions, such as long to double, can lose precision. Integer overflow does not automatically raise an exception. Use modest values in beginner arithmetic.\par Syllabus reading: Liang Ch 2. CLO-1.}
\end{frame}

\begin{frame}[fragile,t]{Quotient and remainder decompose a quantity}
\begin{itemize}
\item Whole-unit conversion uses integer division to count complete larger units.
\item Remainder keeps the part that has not yet been converted.
\item Convert the remainder again when more than two units are needed.
\end{itemize}
\vspace{1mm}\begin{block}{Example}
\begin{lstlisting}
3672 seconds
1 complete hour leaves 72 seconds
1 complete minute leaves 12 seconds
\end{lstlisting}
\end{block}
\note{Explain the mechanism using the displayed example. Whole-unit conversion uses integer division to count complete larger units. Remainder keeps the part that has not yet been converted. Convert the remainder again when more than two units are needed.\par Syllabus reading: Liang Ch 2. CLO-1.}
\end{frame}

\begin{frame}[fragile,t]{Casting changes a value at a specific point}
\begin{itemize}
\item A cast affects the expression at the location where it occurs.
\item Casting an already truncated quotient cannot restore its discarded fraction.
\item Promote an operand before division when a fractional result is required.
\end{itemize}
\vspace{1mm}\begin{block}{Example}
\begin{lstlisting}
(double) (7 / 2) gives 3.0
(double) 7 / 2 gives 3.5
\end{lstlisting}
\end{block}
\note{Explain the mechanism using the displayed example. A cast affects the expression at the location where it occurs. Casting an already truncated quotient cannot restore its discarded fraction. Promote an operand before division when a fractional result is required.\par Syllabus reading: Liang Ch 2. CLO-1.}
\end{frame}


\begin{frame}[fragile,t]{Arithmetic follows the expression tree}
\begin{center}\begin{tikzpicture}
\node[decision] (a) at (0,0) {/};
\node[decision] (b) at (-2,-1.25) {+};
\node[alt] (c) at (2,-1.25) {2.0};
\node[box] (d) at (-3.6,-2.65) {7};
\node[box] (e) at (-0.3,-2.65) {8};
\draw[arr] (a)--(b);
\draw[arr] (a)--(c);
\draw[arr] (b)--(d);
\draw[arr] (b)--(e);
\end{tikzpicture}\end{center}
\begin{block}{What to notice}The denominator 2.0 makes the final division floating-point.\end{block}
\note{Trace each labelled connection. The denominator 2.0 makes the final division floating-point.}
\end{frame}
\begin{frame}[fragile,t]{Key distinctions}
\small\renewcommand{\arraystretch}{1.5}
\rowcolors{2}{pale}{white}
\begin{tabularx}{\textwidth}{>{\raggedright\arraybackslash}X>{\raggedright\arraybackslash}X>{\raggedright\arraybackslash}X}
\rowcolor{navy}
\color{white}\bfseries Expression & \color{white}\bfseries Operand arithmetic & \color{white}\bfseries Result\\
7 / 2 & Integer division & 3\\
7 / 2.0 & Floating-point division & 3.5\\
7 \% 2 & Integer remainder & 1\\
(int) 7.9 & Narrowing conversion & 7\\
(double) (7 / 2) & Cast after division & 3.0\\
\end{tabularx}
\vspace{3mm}\footnotesize These distinctions determine which operation or control structure fits the problem.
\note{\par Syllabus reading: Liang Ch 2. CLO-1.}
\end{frame}

\begin{frame}[fragile,t]{First example: execution}
\begin{lstlisting}
int total = 145;
int count = 2;
double wrong = total / count;
double correct = (double) total / count;
System.out.println(wrong);
System.out.println(correct);
\end{lstlisting}
\note{This is the main body from Java\_Examples/Lecture05.java. The source file includes the complete class. Predict the result before the next slide.\par Syllabus reading: Liang Ch 2. CLO-1.}
\end{frame}

\begin{frame}[fragile,t]{First example: result}
\begin{columns}[T,onlytextwidth]\begin{column}{0.60\textwidth}
72.0\par 72.5\par The first expression performs integer division before assignment.\par The cast changes an operand before division.
\end{column}\begin{column}{0.35\textwidth}\small
Precedence and grouping

Conversions and numeric limits
\end{column}\end{columns}
\note{Expected output and interpretation: 72.0
72.5
The first expression performs integer division before assignment.
The cast changes an operand before division.\par Syllabus reading: Liang Ch 2. CLO-1.}
\end{frame}

\begin{frame}[fragile,t]{Guided practice}
\begin{columns}[T,onlytextwidth]\begin{column}{0.60\textwidth}
Convert 3672 seconds to hours, minutes and seconds using / and \%.
\end{column}\begin{column}{0.35\textwidth}\small
Attempt the solution before continuing.

Use the definitions and examples from this lecture.
\end{column}\end{columns}
\note{Give students 8 minutes to attempt the problem. The following slides provide a complete solution. Original solution guidance: hours = 3672 / 3600 = 1
minutes = (3672 \% 3600) / 60 = 1
seconds = 3672 \% 60 = 12\par Syllabus reading: Liang Ch 2. CLO-1.}
\end{frame}

\begin{frame}[fragile,t]{Solution strategy}
\begin{columns}[T,onlytextwidth]\begin{column}{0.60\textwidth}
\begin{itemize}
\item Find complete hours with seconds / 3600.
\item Keep seconds \% 3600 for the part after complete hours.
\item Divide that remainder by 60, then use \% 60 for remaining seconds.
\end{itemize}
\end{column}\begin{column}{0.35\textwidth}\small
The next program uses fixed inputs so its result can be reproduced.
\end{column}\end{columns}
\note{Discuss why each step is needed. State the input assumptions before executing the program.\par Syllabus reading: Liang Ch 2. CLO-1.}
\end{frame}

\begin{frame}[fragile,t]{Complete solution}
\begin{lstlisting}
public class Solution05 {
  public static void main(String[] args) throws Exception {
    int totalSeconds = 3672;
    int hours = totalSeconds / 3600;
    int remainder = totalSeconds % 3600;
    int minutes = remainder / 60;
    int seconds = remainder % 60;
    System.out.println(hours + ":" + minutes + ":" + seconds);
  }
}
\end{lstlisting}
\note{Complete runnable file: Java\_Examples/Solution05.java. Inputs are fixed in the program.  \par Syllabus reading: Liang Ch 2. CLO-1.}
\end{frame}

\begin{frame}[fragile,t]{Execution trace}
\small\renewcommand{\arraystretch}{1.5}
\rowcolors{2}{pale}{white}
\begin{tabularx}{\textwidth}{>{\raggedright\arraybackslash}X>{\raggedright\arraybackslash}X>{\raggedright\arraybackslash}X}
\rowcolor{navy}
\color{white}\bfseries Variable & \color{white}\bfseries Calculation & \color{white}\bfseries Value\\
hours & 3672 / 3600 & 1\\
remainder & 3672 \% 3600 & 72\\
minutes & 72 / 60 & 1\\
seconds & 72 \% 60 & 12\\
\end{tabularx}
\vspace{3mm}\footnotesize Follow the changing state in execution order.
\note{Manually derived trace for the complete solution. Find complete hours with seconds / 3600. Keep seconds \% 3600 for the part after complete hours. Divide that remainder by 60, then use \% 60 for remaining seconds.\par Syllabus reading: Liang Ch 2. CLO-1.}
\end{frame}

\begin{frame}[fragile,t]{Solution result}
\begin{columns}[T,onlytextwidth]\begin{column}{0.60\textwidth}
1:1:12
\end{column}\begin{column}{0.35\textwidth}\small
Why this result follows

Divide that remainder by 60, then use \% 60 for remaining seconds.
\end{column}\end{columns}
\note{Expected result: 1:1:12\par Syllabus reading: Liang Ch 2. CLO-1.}
\end{frame}

\begin{frame}[fragile,t]{A common incorrect approach}
double average = (double) (7 / 2);
\vspace{1mm}\begin{block}{Example}
\begin{lstlisting}
The cast is too late. Use (double) 7 / 2 to obtain 3.5.
\end{lstlisting}
\end{block}
\note{Ask students to predict the failure before discussing the explanation. The right side gives the correction.\par Syllabus reading: Liang Ch 2. CLO-1.}
\end{frame}

\begin{frame}[fragile,t]{Boundary and failure checks}
\begin{columns}[T,onlytextwidth]\begin{column}{0.60\textwidth}
Check the stated input assumptions.

Write a temperature converter. Test 0 and 100 Celsius. Explain why 9 / 5 gives the wrong conversion factor.
\end{column}\begin{column}{0.35\textwidth}\small
hours = 3672 / 3600 = 1\par minutes = (3672 \% 3600) / 60 = 1\par seconds = 3672 \% 60 = 12
\end{column}\end{columns}
\note{Use the specification to derive each expected result. Exercise solution: hours = 3672 / 3600 = 1
minutes = (3672 \% 3600) / 60 = 1
seconds = 3672 \% 60 = 12\par Syllabus reading: Liang Ch 2. CLO-1.}
\end{frame}

\begin{frame}[fragile,t]{Check your understanding}
\begin{columns}[T,onlytextwidth]\begin{column}{0.60\textwidth}
Evaluate 10 - 3 * 2 and (10 - 3) * 2. What does (int) -2.9 produce?
\end{column}\begin{column}{0.35\textwidth}\small
Write a prediction and explain the rule that supports it.
\end{column}\end{columns}
\note{Answer: 4, 14 and -2. The cast truncates toward zero.\par Syllabus reading: Liang Ch 2. CLO-1.}
\end{frame}

\begin{frame}[fragile,t]{Answer and explanation}
\begin{columns}[T,onlytextwidth]\begin{column}{0.60\textwidth}
4, 14 and -2. The cast truncates toward zero.
\end{column}\begin{column}{0.35\textwidth}\small
The cast is too late. Use (double) 7 / 2 to obtain 3.5.
\end{column}\end{columns}
\note{Reconnect the answer to the relevant definition rather than treating it as a fact to memorise.\par Syllabus reading: Liang Ch 2. CLO-1.}
\end{frame}


\begin{frame}[fragile,t]{Compare cases and predict outcomes}
\small\renewcommand{\arraystretch}{1.5}\rowcolors{2}{pale}{white}
\begin{tabularx}{\textwidth}{>{\raggedright\arraybackslash}X>{\raggedright\arraybackslash}X>{\raggedright\arraybackslash}X}\rowcolor{navy}
\color{white}\bfseries Expression & \color{white}\bfseries Result & \color{white}\bfseries Reason\\
(7 + 8) / 2 & 7 & Integer quotient\\
(7 + 8) / 2.0 & 7.5 & Floating-point divisor\\
(double) ((7 + 8) / 2) & 7.0 & The fraction was already lost\\
\end{tabularx}\par\vspace{3mm}\begin{exampleblock}{Discuss}Explain each expected result before running the example. Which case would reveal a wrong assumption?\end{exampleblock}
\note{Read the first two columns and invite predictions before discussing the outcome.}
\end{frame}

\begin{frame}[fragile,t]{Apply the idea}
\begin{alertblock}{Independent practice}Convert 5,439 seconds to hours, minutes and seconds. Explain why quotient and remainder are both needed.\end{alertblock}\vspace{3mm}\begin{enumerate}\item Work through the input by hand.\item Write or correct the program.\item Justify the expected result.\end{enumerate}\note{Allow 3–5 minutes, then compare answers in pairs. Convert 5,439 seconds to hours, minutes and seconds. Explain why quotient and remainder are both needed.}
\end{frame}

\begin{frame}[fragile,t]{Solution and reasoning}
\begin{lstlisting}
int seconds = 5439;
int hours = seconds / 3600;
int remainder = seconds % 3600;
int minutes = remainder / 60;
int leftover = remainder % 60;
System.out.println(hours + ":" + minutes + ":" + leftover);
\end{lstlisting}\vspace{1mm}\small\begin{itemize}
\item The first quotient gives 1 complete hour; 1,839 seconds remain.
\item The next quotient gives 30 minutes; the remainder is 39 seconds.
\item Output: 1:30:39.
\end{itemize}\note{Answer: The first quotient gives 1 complete hour; 1,839 seconds remain. The next quotient gives 30 minutes; the remainder is 39 seconds. Output: 1:30:39.}
\end{frame}

\begin{frame}[fragile,t]{Division and compound assignment}
\begin{itemize}
\item Integer division discards the fractional part; a floating{-}point operand changes the arithmetic.
\item A compound assignment includes conversion back to the left variable type.
\item Therefore x /= 3.0 for an int x is not equivalent to the uncompilable x = x / 3.0.
\end{itemize}
\vspace{3mm}\footnotesize\textcolor{accent}{Source: supplied ISB campus slides. See speaker notes and ISB\_Source\_Map.md.}
\note{Adapted from the supplied ISB campus folder: Lecture{-}4.pptx, slides 13, 16, 17. Clarified the implicit conversion in compound assignment, which the source equivalence omits.}
\end{frame}

\begin{frame}[fragile,t]{Worked problem: Division and compound assignment}
\begin{block}{Task}Compare 5/2, 5.0/2 and int x=10 followed by x/=3.0. Explain each result.\end{block}
\small Input: Values are fixed in the program for a reproducible demonstration.\par\vspace{3mm}\begin{exampleblock}{Before running}Predict the output and identify the variables that change.\end{exampleblock}
\note{Adapted from the supplied ISB campus folder: Lecture{-}4.pptx, slides 13, 16, 17. Clarified the implicit conversion in compound assignment, which the source equivalence omits.}
\end{frame}

\begin{frame}[fragile,t]{Complete ISB example (1/1)}
\begin{lstlisting}
public class IsbLecture05 {
  public static void main(String[] args) throws Exception {
    System.out.println(5 / 2);
    System.out.println(5.0 / 2);
    int x = 10;
    x /= 3.0;
    System.out.println(x);
  }

}
\end{lstlisting}
\note{Adapted from the supplied ISB campus folder: Lecture{-}4.pptx, slides 13, 16, 17. Clarified the implicit conversion in compound assignment, which the source equivalence omits. Complete file: ISB\_Java\_Examples/IsbLecture05.java.}
\end{frame}

\begin{frame}[fragile,t]{Worked trace and expected result}
\small\renewcommand{\arraystretch}{1.4}\rowcolors{2}{pale}{white}
\begin{tabularx}{\textwidth}{>{\raggedright\arraybackslash}X>{\raggedright\arraybackslash}X>{\raggedright\arraybackslash}X}\rowcolor{navy}
\color{white}\bfseries Expression & \color{white}\bfseries Arithmetic before storage & \color{white}\bfseries Result\\
5 / 2 & Integer & 2\\
5.0 / 2 & double & 2.5\\
x /= 3.0, x initially 10 & double then int conversion & 3\\
\end{tabularx}\par\vspace{2mm}
\begin{block}{Expected console output}\ttfamily\footnotesize 2\par 2.5\par 3\end{block}
\note{Adapted from the supplied ISB campus folder: Lecture{-}4.pptx, slides 13, 16, 17. Clarified the implicit conversion in compound assignment, which the source equivalence omits. Trace and expected values independently worked for this edition.}
\end{frame}

\begin{frame}[fragile,t]{Extend the ISB example}
\begin{alertblock}{Practice}Rewrite x /= 3.0 with an explicit assignment that compiles and preserves the demonstrated conversion.\end{alertblock}\vspace{3mm}\begin{exampleblock}{Explain your reasoning}x = (int) (x / 3.0); The cast occurs after division.\end{exampleblock}
\note{Adapted from the supplied ISB campus folder: Lecture{-}4.pptx, slides 13, 16, 17. Clarified the implicit conversion in compound assignment, which the source equivalence omits. Answer: x = (int) (x / 3.0); The cast occurs after division.}
\end{frame}
\begin{frame}[fragile,t]{Lecture summary}
\begin{columns}[T,onlytextwidth]\begin{column}{0.60\textwidth}
\begin{itemize}
\item suitable primitive types
\item precedence and integer division
\item casts and conversions deliberately
\end{itemize}
\end{column}\begin{column}{0.35\textwidth}\small
\begin{itemize}
\item Find complete hours with seconds / 3600.
\item Keep seconds \% 3600 for the part after complete hours.
\item Divide that remainder by 60, then use \% 60 for remaining seconds.
\end{itemize}
\end{column}\end{columns}
\note{Ask students to give one concrete example for the concepts listed. Use the worked solution to connect the topics.\par Syllabus reading: Liang Ch 2. CLO-1.}
\end{frame}

\begin{frame}[fragile,t]{After-class practice}
\begin{columns}[T,onlytextwidth]\begin{column}{0.60\textwidth}
Write a temperature converter. Test 0 and 100 Celsius. Explain why 9 / 5 gives the wrong conversion factor.
\end{column}\begin{column}{0.35\textwidth}\small
Syllabus reading\par Liang Ch 2

Next lecture: Updates and formatted output
\end{column}\end{columns}
\note{Reading labels follow the supplied outline. Check topic headings against the available textbook edition. Require a program and expected results for the practice task.\par Syllabus reading: Liang Ch 2. CLO-1.}
\end{frame}
\end{document}
